\chapter{Categorical and graphical affine algebra}\label{sec:appendix-gla}

This appendix develops the categorical background used by the graphical
calculi in this thesis.  Its purpose is not to survey category theory in full,
but to make the passage from ordinary composition to Graphical Affine Algebra
(GAA) explicit and self-contained.  The central idea is simple: a diagram is a
morphism, joining diagrams eliminates their common variables, and placing
diagrams side by side forms independent systems.  Category theory isolates
these two operations; symmetric monoidal categories make them graphical; and
GAA specializes the resulting language to systems of affine equations.

The development proceeds in four stages.  First we recall categories, functors,
and the category of relations.  We then add a symmetric monoidal product and
explain the corresponding string-diagram notation.  Next we introduce PROPs,
monoid and Frobenius structures, which provide the algebraic operations for
copying, deleting, adding, and constraining variables.  Finally we define the
PROP of affine relations and state the presentation theorem identifying it
with GAA.

\section{Categories: systems and their composition}

\begin{definition}[Category]
A \emph{category} $\mathcal C$ consists of:
\begin{itemize}
    \item a class of objects $A,B,\ldots$;
    \item for each pair of objects, a class $\mathcal C(A,B)$ of morphisms
          $f:A\to B$;
    \item for every $A$, an identity morphism $\mathrm{id}_A:A\to A$;
    \item for $f:A\to B$ and $g:B\to C$, a composite
          $g\circ f:A\to C$.
\end{itemize}
Composition is associative and identities are neutral:
\[
 h\circ(g\circ f)=(h\circ g)\circ f,
 \qquad
 \mathrm{id}_B\circ f=f=f\circ\mathrm{id}_A.
\]
\end{definition}

A category therefore retains only what is required to connect processes.  An
object records the type of an interface, while a morphism records a process
with a specified input and output.  The equation $g\circ f$ is meaningful
exactly when the output type of $f$ agrees with the input type of $g$.
Associativity says that a composite process does not depend on how it is
parenthesized.

Two examples will be used repeatedly.

\begin{example}[Functions]
The category $\mathbf{Set}$ has sets as objects and functions as morphisms.
Composition and identities are the usual ones.  A function is deterministic:
each input has exactly one output.
\end{example}

\begin{example}[Relations]\label{ex:appendix-rel}
The category $\mathbf{Rel}$ has sets as objects and binary relations
$R\subseteq X\times Y$ as morphisms $R:X\to Y$.  Given $R:X\to Y$ and
$S:Y\to Z$, their composite is
\[
 S\circ R
 =\{(x,z)\mid \text{there exists }y\in Y
              \text{ with }xRy\text{ and }ySz\}.
\]
The identity on $X$ is the equality relation
$\{(x,x)\mid x\in X\}$.  Thus relational composition eliminates the shared
variable $y$.  Unlike a function, a relation may associate an input with no
output or with several outputs; this is why relations, rather than functions,
are the appropriate semantics for systems of equations and inequalities.
\end{example}

\begin{definition}[Functor]
A \emph{functor} $F:\mathcal C\to\mathcal D$ assigns an object $FA$ to every
object $A$ and a morphism $Ff:FA\to FB$ to every morphism $f:A\to B$, while
preserving identities and composition:
\[
 F(\mathrm{id}_A)=\mathrm{id}_{FA},
 \qquad
 F(g\circ f)=Fg\circ Ff.
\]
It is \emph{faithful} when every map
$\mathcal C(A,B)\to\mathcal D(FA,FB)$ is injective, and \emph{full} when every
such map is surjective.
\end{definition}

In this work, a functor usually interprets syntax as semantics.  Faithfulness
means that semantically equal diagrams are already equal in the calculus;
fullness means that every semantic process of the intended kind can be drawn.
These are the categorical forms of completeness and expressivity,
respectively.

\section{Parallel composition and string diagrams}

Ordinary categorical composition describes processes in sequence.  To describe
independent processes in parallel we need a second operation.

\begin{definition}[Symmetric monoidal category]\label{appendix:smc}
A \emph{symmetric monoidal category} is a category $\mathcal C$ equipped with a
bifunctor
\[
 \otimes:\mathcal C\times\mathcal C\longrightarrow\mathcal C,
\]
a monoidal unit $I$, and coherent natural isomorphisms for associativity,
units, and symmetry.  In particular, for morphisms $f:A\to B$ and $g:C\to D$
there is a morphism
\[
 f\otimes g:A\otimes C\longrightarrow B\otimes D,
\]
and the interchange law holds:
\[
 (f_2\circ f_1)\otimes(g_2\circ g_1)
 = (f_2\otimes g_2)\circ(f_1\otimes g_1).
\]
The symmetry $\sigma_{A,B}:A\otimes B\to B\otimes A$ exchanges the two
components.
\end{definition}

Mac Lane's coherence theorem allows the associativity and unit isomorphisms to
be suppressed without ambiguity~\cite{maclane1998categories}.  We use this
strict notation throughout.

A symmetric monoidal category has a sound two-dimensional notation: its
\emph{string diagrams}~\cite{joyal1991geometry,selinger2010survey}.  A wire
represents an object and a box with input and output wires represents a
morphism.  Reading from left to right, joining the outputs of $f$ to the inputs
of $g$ represents $g\circ f$; placing $f$ above $g$ represents $f\otimes g$.
A bare wire is an identity and crossing wires represents symmetry.  The
interchange law is built into the geometry: disjoint boxes may move past one
another because their relative horizontal order carries no information.
Consequently, diagrams are considered equal up to deformations that preserve
connectivity, boundary, and the order of ports on each box.  The basic
translation between categorical notation and string diagrams is displayed in
Figure~\ref{fig:appendix-string-diagram-dictionary}.

\begin{figure}[H]
\centering
\begin{tikzpicture}[axpic]
  % Identity and a general morphism
  \draw (-0.6,1.1) -- (0.6,1.1);
  \node at (0,0.45) {$\mathrm{id}_A:A\to A$};

  \pic[val=f] (single) at (2.6,1.1) {diagram};
  \node at (2.6,0.45) {$f:A\to B$};

  % Sequential composition
  \pic[val=f] (cf) at (4.6,1.1) {diagram};
  \pic[val=g] (cg) at (5.9,1.1) {diagram};
  \draw (cf-out-0) -- (cg-in-0);
  \node at (5.25,0.45) {$g\circ f$};

  % Tensor product
  \pic[val=f] (tf) at (7.8,1.65) {diagram};
  \pic[val=g] (tg) at (7.8,0.75) {diagram};
  \node at (7.8,0.1) {$f\otimes g$};

  % Symmetry
    \draw (9.5,1.45) to[out=0,in=180] (10.7,0.75);
    \draw (9.5,0.75) to[out=0,in=180] (10.7,1.45);
  \node at (10.1,0.1) {$\sigma_{A,B}$};
\end{tikzpicture}
\caption{String-diagram notation for the structure of a symmetric monoidal
category.  Diagrams are read from left to right: horizontal connection is
composition, while vertical juxtaposition is the tensor product.}
\label{fig:appendix-string-diagram-dictionary}
\end{figure}

The labels beneath the pictures record the corresponding morphisms.  In
particular, the two wires in the drawing of $f\otimes g$ do not interact: they
represent processes performed in parallel.  The crossing represents the
canonical reordering $A\otimes B\to B\otimes A$, not an additional operation
chosen for the particular objects.

\begin{figure}[H]
    \centering
\begin{tikzpicture}[axpic]
  \pic[val=f] (f) at (0,0.7) {diagram};
  \pic[val=g] (g) at (1.5,0.7) {diagram};
  \draw (f-out-0) -- (g-in-0);
  \pic[val=h] (h) at (0,-0.7) {diagram};
  \pic[val=k] (k) at (1.5,-0.7) {diagram};
  \draw (h-out-0) -- (k-in-0);
  \node at (3.0,0) {$=$};
  \pic[val=f] (f2) at (4.0,0.7) {diagram};
  \pic[val=g] (g2) at (5.5,0.7) {diagram};
  \draw (f2-out-0) -- (g2-in-0);
  \pic[val=h] (h2) at (4.0,-0.7) {diagram};
  \pic[val=k] (k2) at (5.5,-0.7) {diagram};
  \draw (h2-out-0) -- (k2-in-0);
  \node[below=2 of f-node] {$(g\circ f)\otimes(k\circ h)$};
  \node[below=2 of f2-node] {$(g\otimes k)\circ(f\otimes h)$};
\end{tikzpicture}
\caption{The interchange law expressed as a string diagram.}
\label{fig:appendix-interchange}
\end{figure}

For $\mathbf{Rel}$ the monoidal product is Cartesian product on objects.  If
$R:X\to Y$ and $S:U\to V$, then
\[
 R\otimes S:X\times U\longrightarrow Y\times V,
 \qquad
 ((x,u),(y,v))\in R\otimes S
 \Longleftrightarrow xRy\text{ and }uSv.
\]
The singleton set is the unit.  Sequential composition eliminates shared
variables; monoidal composition juxtaposes independent systems.

\section{PROPs and graphical theories}

The graphical calculi used in this thesis have one generating wire type.  An
interface is therefore determined only by its number of wires.

\begin{definition}[PROP]
A \emph{PROP} is a strict symmetric monoidal category whose objects are the
natural numbers, whose tensor product on objects is addition, and whose
monoidal unit is $0$.  A morphism $m\to n$ is drawn with $m$ input wires and
$n$ output wires.
\end{definition}

Thus $1$ is the basic wire, $n=1^{\otimes n}$ is a bundle of $n$ wires, and
$m\otimes n=m+n$.  Matrices provide the standard example: for a field
$\mathbb K$, let $\mathbf{Mat}_{\mathbb K}$ have morphisms $m\to n$ given by
$n\times m$ matrices.  Composition is matrix multiplication, tensor is block
diagonal sum, and the symmetry is a permutation matrix.

\section{Monoids, comonoids, and Frobenius structure}

The algebraic nodes in GAA are not arbitrary operations.  They arise from
standard algebraic structures internal to a symmetric monoidal category.

\begin{definition}[Monoid and comonoid]
A \emph{commutative monoid} on an object $A$ consists of
$\mu:A\otimes A\to A$ and $\eta:I\to A$ satisfying associativity, unitality,
and commutativity.  Dually, a \emph{cocommutative comonoid} consists of
$\delta:A\to A\otimes A$ and $\varepsilon:A\to I$ satisfying coassociativity,
counitality, and cocommutativity.
\end{definition}

On a vector space, the intended monoid is addition:
$\mu(x_1,x_2)=x_1+x_2$ and $\eta(*)=0$.  The intended comonoid copies and
discards data:
$\delta(x)=(x,x)$ and $\varepsilon(x)=*$.  Copying is linear, but its converse
is not a function; this is the first indication that linear relations are the
natural ambient category.

\begin{definition}[Special commutative Frobenius algebra]
A commutative monoid $(\mu,\eta)$ and a cocommutative comonoid
$(\delta,\varepsilon)$ on the same object form a \emph{special commutative
Frobenius algebra} when
\[
 (\mu\otimes\mathrm{id})\circ(\mathrm{id}\otimes\delta)
 =\delta\circ\mu
 =(\mathrm{id}\otimes\mu)\circ(\delta\otimes\mathrm{id})
\]
and $\mu\circ\delta=\mathrm{id}$.
\end{definition}

The Frobenius and special laws imply the \emph{spider theorem}: every connected
diagram built only from $\mu,\eta,\delta,\varepsilon$ is determined by its
number of input and output wires.  It may therefore be contracted to one
multi-input, multi-output node.  This is the algebraic reason that wires can be
split, merged, bent, and reconnected without recording an arbitrary parse
tree.

The cup $\delta\circ\eta:I\to A\otimes A$ and cap
$\varepsilon\circ\mu:A\otimes A\to I$ induced by a special commutative
Frobenius algebra satisfy the snake equations.  They make $A$ self-dual and allow an
input wire to be bent into an output wire, or conversely.  A symmetric monoidal
category equipped coherently with such a structure on each object is called a
\emph{hypergraph category}~\cite{fong2019hypergraph}.  This compact structure
is what makes the converse of a relation available graphically by reflection.

\section{The PROP of affine relations}

Fix a field $\mathbb K$; all applications in this thesis use
$\mathbb K=\mathbb R$.

\begin{definition}[Affine relation]
An \emph{affine relation} $R:m\to n$ is an affine subspace (or the empty set)
\[
 R\subseteq\mathbb K^m\oplus\mathbb K^n.
\]
We write $x\mathrel{R}y$ when $(x,y)\in R$.  The PROP
$\mathbf{AffRel}_{\mathbb K}$ has natural numbers as objects and affine
relations as morphisms.  Its identity, composition, and tensor are
\begin{align*}
 \mathrm{id}_n
   &=\{(x,x)\mid x\in\mathbb K^n\},\\
 S\circ R
   &=\{(x,z)\mid \exists y\in\mathbb K^n:
                    x\mathrel{R}y\text{ and }y\mathrel{S}z\},\\
 R\otimes S
   &=\{((x,u),(y,v))\mid x\mathrel{R}y,
                              \ u\mathrel{S}v\}.
\end{align*}
After the canonical identification
$\mathbb K^m\oplus\mathbb K^p\cong\mathbb K^{m+p}$, the tensor on objects is
addition.
\end{definition}

The composite is again an affine subspace or empty: the simultaneous
constraints $(x,y)\in R$ and $(y,z)\in S$ form an affine system, and projecting
away $y$ preserves affinity.  Hence composition of affine relations is exactly
affine variable elimination.

Every matrix $A\in\mathbb K^{n\times m}$ determines the functional linear
relation
\[
 \operatorname{graph}(A)
 =\{(x,y)\in\mathbb K^m\oplus\mathbb K^n\mid y=Ax\}.
\]
Matrix multiplication is recovered by relational composition:
\[
 \operatorname{graph}(B)\circ\operatorname{graph}(A)
 =\operatorname{graph}(BA).
\]
Affine relations are nevertheless strictly more general than affine maps.
Linear maps embed as the special case with zero translation.  For example, the
converse
\[
 \operatorname{graph}(A)^{\mathrm{op}}
 =\{(y,x)\mid y=Ax\}
\]
may be partial or multivalued.  More generally, every affine relation admits a presentation by affine linear
equations:
\[
 R=\{(x,y)\mid Ex+b=Fy\}
\]
for suitable matrices $E,F$ and vector $b$.  This affine cospan form is the
semantic origin of GAA normal forms.

\section{Graphical Affine Algebra}

Graphical Affine Algebra is a graphical presentation of
$\mathbf{AffRel}_{\mathbb K}$~\cite{bonchi2019affine}.  Its diagrams are
generated by two interacting families of nodes, one scalar box for each
$k\in\mathbb K$, and the affine point $\rtop:0\to1$.  We read diagrams from
left to right.  The following table gives the type and exact relation denoted
by every basic orientation; the reversed orientation denotes the converse
relation.

\begin{table}[H]
\centering
\caption{Generators of $\mathbf{GAA}$ and their affine-relation semantics.}
\label{tab:appendix-gla-generators}
\begin{tabular}{c c p{8.1cm}}
\\
\toprule
\textbf{diagram} & \textbf{type} & \textbf{interpretation} \\
\midrule
\begin{tikzpicture}[axpic]\pic (g) at (0,0) {multiplication};\end{tikzpicture}
& $2\to1$ & $\{((x_1,x_2),y)\mid y=x_1+x_2\}$ \\
\begin{tikzpicture}[axpic]\pic (g) at (0,0) {zero};\end{tikzpicture}
& $0\to1$ & $\{(*,0)\}$ \\
\begin{tikzpicture}[axpic]\pic (g) at (0,0) {copy};\end{tikzpicture}
& $1\to2$ & $\{(x,(y_1,y_2))\mid y_1=x=y_2\}$ \\
\begin{tikzpicture}[axpic]\pic (g) at (0,0) {discard};\end{tikzpicture}
& $1\to0$ & $\{(x,*)\mid x\in\mathbb K\}$ \\
\begin{tikzpicture}[axpic]\pic[val=k] (g) at (0,0) {scalar};\end{tikzpicture}
& $1\to1$ & $\{(x,y)\mid y=kx\}$, for $k\in\mathbb K$ \\
\begin{tikzpicture}[axpic]\pic (g) at (0,0) {affine};\end{tikzpicture}
& $0\to1$ & $\{(*,1)\}$ \\
\bottomrule
\end{tabular}
\end{table}

The white node is addition and the black node is equality.  Their reflected
forms have the following readings:
\begin{itemize}
    \item white comultiplication $1\to2$ imposes $x=y_1+y_2$;
    \item white counit $1\to0$ imposes $x=0$;
    \item black multiplication $2\to1$ imposes $x_1=x_2=y$;
    \item black unit $0\to1$ allows an arbitrary value of $y$;
    \item a reflected scalar imposes $x=ky$.
\end{itemize}
Thus direction distinguishes solving an equation forward from treating the
same equation as a constraint.  When $k\ne0$, a reflected $k$-box is equivalent
to a forward $k^{-1}$-box; when $k=0$, its semantics is the non-functional
constraint $x=0$ with arbitrary output.

The equational theory of GAA can be organized into five layers.

\begin{enumerate}[label=\textnormal{(\roman*)}]
    \item The addition node and zero form a commutative monoid, while copy and
    discard form a cocommutative comonoid.  Reflecting these equations gives
    the corresponding structures in the opposite direction.

    \item Each color, together with its reflection, is a special commutative
    Frobenius algebra.  The spider theorem therefore applies separately to
    connected monochromatic regions.

    \item The two colors interact by the bialgebra and Hopf laws.  These laws
    express that copying preserves addition and zero, and that additive
    inverses cancel.  The scalar $-1$ is the antipode.

    \item Scalar boxes obey the arithmetic of $\mathbb K$: consecutive boxes
    multiply, parallel copied-and-added boxes add, $1$ is the identity, and
    scalars commute with the appropriate copying and addition nodes.

    \item The affine point $\rtop$ denotes the constant $1$.  It is copied by
    the black comultiplication and discarded by the black counit; these are the
    affine equations that allow constants and translations to propagate
    through a diagram.
\end{enumerate}

Every equation is sound because both sides denote the same affine subspace.
For example, copying after addition and adding copied components both denote
\[
 \{((x_1,x_2),(y_1,y_2))\mid
       y_1=x_1+x_2=y_2\}.
\]
The equations do more than preserve semantics: they characterize it.

\begin{theorem}[Presentation theorem for GAA]\label{thm:appendix-gla-presentation}
For a field $\mathbb K$, the standard interpretation extends to an isomorphism
of PROPs
\[
 \llbracket{-}\rrbracket:
 \mathbf{GAA}_{\mathbb K}\xrightarrow{\ \cong\ }
 \mathbf{AffRel}_{\mathbb K}.
\]
Consequently, GAA is sound, complete, and expressive for affine
relations~\cite{bonchi2019affine}.
\end{theorem}

The theorem explains why GAA is useful in the main text.  A diagram is not
merely a picture of a chosen matrix calculation: it denotes the entire
solution relation of an affine system.  Sequential composition existentially
eliminates internal variables, tensor combines independent systems, and
reflection exchanges a relation with its converse.  Diagram rewriting is
therefore a complete algebra of affine systems.

A matrix box $A:m\to n$ abbreviates any GAA diagram whose interpretation is
$y=Ax$.  Its reflected box denotes the converse relation.  For compatible
matrices, connecting boxes gives matrix multiplication, while copying inputs,
applying two matrices, and adding their outputs gives the corresponding block
matrix.  A general relation box may be replaced by its cospan form
\[
 x\mathrel{R}y
 \quad\Longleftrightarrow\quad
 Ex=Fy.
\]
Cups and caps bend wires, transposing a diagram from the functional fragment
into its relational converse.  Finally, connected black Frobenius diagrams
contract by the spider theorem.

The particular derived equalities invoked by the proofs in
Chapters~\ref{sec:gpa} and~\ref{sec:glp} are collected below.  They are grouped
as (G1) matrix calculus, (G2) transposition and cospan form, (G3) bending, and
(G4) the black spider laws.  All follow from the linear fragment of the GAA equations summarized above and
are included here so that later diagrammatic derivations can be checked without
expanding every matrix box into scalar generators.

% Rules of the underlying linear and affine algebras (GLA, GAA) that are used
% as background throughout, but are not contributions of this thesis.
\medskip
\noindent
\fbox{\begin{minipage}{0.97\linewidth}
    \begin{center}\textbf{(G1) Matrices}\quad{\footnotesize\cite{bonchi2017interacting,paixao2022highlevel}}\end{center}
    \begin{center}\begin{tabular}{r c c c}
      \diagrow{(comp.)}{
          \pic[val=A] (g1) at (0.250,0.500) {matrix};
          \pic[val=B] (g2) at (2.250,0.500) {matrix};
          \draw (g1-out-0) to[out=0,in=180] (g2-in-0);
      }{
          \pic[val=BA] (g1) at (0.500,0.500) {matrix};
      }
      \diagrow{(row)}{
          \pic[val=A] (g1) at (0.500,1.000) {matrix};
          \pic[val=B] (g2) at (0.500,0.000) {matrix};
          \pic (g3) at (2.050,0.500) {multiplication};
          \draw (g1-out-0) to[out=0,in=180] (g3-in-0);
          \draw (g2-out-0) to[out=0,in=180] (g3-in-1);
      }{
          \pic[val={[A \ B]}] (g1) at (1.000,0.000) {2dmatrix};
          \draw (0.000,0.300) -- (g1-in-0);
          \draw (0.000,-0.300) -- (g1-in-1);
      }
      \diagrow{(antipode)}{
          \pic (g1) at (0.500,0.500) {antipode};
      }{
          \pic[val=-1] (g1) at (0.500,0.500) {scalar};
      }
      \diagrow{(zero sc.)}{
          \pic[val=0] (g1) at (0.500,0.500) {scalar};
      }{
          \pic (g1) at (0.500,0.500) {discard};
          \pic (g2) at (1.750,0.500) {zero};
          \draw (0.000,0.500) -- (g1-in-0);
          \draw (g2-out-0) -- (2.500,0.500);
      }
      \diagrow{(inversion)}{
          \pic[val=k] (g1) at (0.500,0.500) {scalar};
      }{
          \pic[val=k^{-1}] (g1) at (0.500,0.500) {coscalar};
      }
    \end{tabular}\end{center}
\end{minipage}}

\medskip
\noindent
\fbox{\begin{minipage}{0.97\linewidth}
    \begin{center}\textbf{(G2) Transposition and cospan form}\quad{\footnotesize\cite{bonchi2017interacting}}\end{center}
    A diagram may be slid along a cap, exchanging it for its transpose; on
    $\mathbf{GLA}$ diagrams the transpose is the converse relation, so a matrix
    becomes a comatrix.
    \begin{center}\begin{tabular}{r c c c}
      \diagrow{(slide)}{
          \pic[val=A] (g1) at (0.500,1.000) {matrix};
          \draw (g1-out-0) -- (1.600,1.000) arc (90:-90:0.500) -- (0.000,0.000);
      }{
          \pic[val=A] (g1) at (0.500,0.000) {comatrix};
          \draw (g1-out-0) -- (1.600,0.000) arc (-90:90:0.500) -- (0.000,1.000);
      }
      \diagrow{(cospan)}{
          \pic[val=g] (g1) at (0.500,0.500) {relation};
      }{
          \pic[val=E] (g1) at (0.500,0.500) {matrix};
          \pic[val=F] (g2) at (2.250,0.500) {comatrix};
          \draw (g1-out-0) to[out=0,in=180] (g2-in-0);
      }
    \end{tabular}\end{center}
\end{minipage}}

\medskip
\noindent
\fbox{\begin{minipage}{0.97\linewidth}
    \begin{center}\textbf{(G3) Bending}\quad{\footnotesize\cite{bonchi2017interacting,bonchi2019affine}}\end{center}
    The white cap is the black cap with an antipode on one leg; this is what
    turns an equation into a pair of bendable wires.
    \begin{center}\begin{tabular}{r c c c}
      \diagrow{(white cap)}{
          \pic (g1) at (0.500,0.500) {multiplication};
          \pic (g2) at (1.750,0.500) {cozero};
          \draw (g1-out-0) to[out=0,in=180] (g2-in-0);
      }{
          \pic (g1) at (0.500,0.000) {antipode};
          \draw (g1-out-0) -- (1.600,0.000) arc (-90:90:0.500) -- (0.000,1.000);
      }
      \diagrow{(unbend $\oplus$)}{
          \pic (g1) at (0.500,0.500) {comultiplication};
          \pic (g2) at (1.750,0.000) {antipode};
          \draw (g1-out-0) -- (2.500,1.000);
          \draw (g1-out-1) to[out=0,in=180] (g2-in-0);
          \draw (g2-out-0) -- (2.500,0.000);
      }{
          \pic (g1) at (1.000,0.500) {multiplication};
          \draw (0.000,1.000) to[out=0,in=180] (g1-in-0);
          \draw (g1-out-0) -- (1.900,0.500) arc (90:-90:0.500) -- (0.000,-0.500);
      }
    \end{tabular}\end{center}
\end{minipage}}

\medskip
\noindent
\fbox{\begin{minipage}{0.97\linewidth}
    \begin{center}\textbf{(G4) The black structure is a spider}\quad{\footnotesize\cite{bonchi2017interacting}}\end{center}
    Any connected network of black nodes with $a$ inputs and $b$ outputs equals
    the single node $\mathsf{cc}_a;\mathsf{cp}_b$; the two laws below are the
    cases used here.
    \begin{center}\begin{tabular}{r c c c}
      \diagrow{(special)}{
          \pic (g1) at (0.500,0.500) {copy};
          \pic (g2) at (1.750,0.500) {cocopy};
          \draw (g1-out-0) to[out=0,in=180] (g2-in-0);
          \draw (g1-out-1) to[out=0,in=180] (g2-in-1);
      }{
          \draw (0.000,0.500) -- (1.000,0.500);
      }
      \diagrow{(merge/del)}{
          \pic (g1) at (0.500,0.500) {cocopy};
          \pic (g2) at (1.750,0.500) {discard};
          \draw (g1-out-0) -- (g2-in-0);
      }{
          \pic (g1) at (0.500,1.000) {discard};
          \pic (g2) at (0.500,0.000) {discard};
          \draw (0.000,1.000) -- (g1-in-0);
          \draw (0.000,0.000) -- (g2-in-0);
      }
      \diagrow{(split/un)}{
          \pic (g1) at (0.500,0.500) {codiscard};
          \pic (g2) at (1.750,0.500) {copy};
          \draw (g1-out-0) -- (g2-in-0);
      }{
          \pic (g1) at (0.500,1.000) {codiscard};
          \pic (g2) at (0.500,0.000) {codiscard};
          \draw (g1-out-0) -- (1.250,1.000);
          \draw (g2-out-0) -- (1.250,0.000);
      }
    \end{tabular}\end{center}
\end{minipage}}


The affine point additionally turns homogeneous matrix equations into affine
ones: together with copying and scalar multiplication it constructs any
constant vector $b$, and hence any constraint $Ex+b=Fy$.  The later calculi in
this thesis extend this affine base with order and optimization generators.
Their new equations are studied in the main chapters, whereas the categorical
and affine rules on which they rely are exactly those established in this
appendix.
